% =====================================================================
% Supplemental tables for the Radiology: AI submission.
%
% E1  Three-way classification of the non-image variables (FINDING 5).
%     The manuscript's Table 4 carries only the G / A / C key; this is the
%     per-variable detail behind it.
% E2  Provenance validation of the hemorrhage benchmark's slice ordering
%     and subject identifier (FINDING 8), each test shown against the value
%     the same statistic returns when the structure being tested is
%     destroyed. The manuscript's Table 2 note states the same statistics in
%     prose; this is the tabular form with the destroyed-structure column.
%
% The benchmark-eligibility record that used to be Table E1 here is now
% Table 1 OF THE MANUSCRIPT, in full, and is not repeated in this file.
%
% ANONYMIZATION: no author, institution, repository, persistent identifier,
% software name or file path appears below. Data-contributing centers are
% referred to by count, not by name.
%
% EXCLUDED VOCABULARY: this file is deliberately free of the token list the
% submission bans; that list is kept in the working notes, not here, so this
% file passes its own check.
%
% Requires: booktabs, siunitx (\num), threeparttable optional.
% =====================================================================

% ---------------------------------------------------------------- Table E1
\begin{table*}[t]
\centering
\footnotesize
\caption{Non-image variables available to the zero-image models, classified three ways;
the manuscript's Table 4 carries only the one-letter key. \textbf{G}, geometry and acquisition: where in the stack the slice
sits and how the image was made. \textbf{A}, administrative: how the data were
catalogued or distributed. \textbf{C}, legitimate clinical non-image predictor.
Values are single-variable area under the receiver operating characteristic
curve at the patient level on each arm's own held-out data, so they show what
each field carries on its own. The headline result of this study rests on class
G alone.}
\label{tab:taxonomy}
\begin{tabular}{@{}p{0.19\textwidth} p{0.29\textwidth} p{0.15\textwidth} p{0.23\textwidth}@{}}
\toprule
Benchmark arm & G, geometry and acquisition & A, administrative & C, clinical \\
\midrule
RSNA ICH (6 label columns) & relative position; series depth 0.601; reconstruction intercept 0.553; imaging plane 0.500; reconstruction slope 0.500 & \emph{none in the file} & \emph{none in the file} \\
\addlinespace[2pt]
fastMRI Prostate, T2 & relative position; volume depth 0.542 & distribution-archive bucket 0.596 & \emph{none in the file} \\
fastMRI Prostate, DWI & relative position; volume depth 0.481 & distribution-archive bucket 0.407 & \emph{none in the file} \\
\addlinespace[2pt]
DeepLesion (8 arms) & published normalized $z$; display window setting 0.518--0.938; scanned range 0.520--0.733; section thickness 0.533--0.694; in-plane spacing 0.534--0.626; matrix size 0.499--0.502; volume depth 0.374--0.705 & \emph{none in the file} & age 0.473--0.566; sex 0.449--0.562 \\
\addlinespace[2pt]
fastMRI+ knee, 2 arms & relative position; readout columns 0.494--0.578; sequence 0.469--0.573; volume depth 0.549--0.553 & \emph{none in the file} & \emph{none in the file} \\
\addlinespace[2pt]
Duke Breast & relative position; volume depth; manufacturer; scanner model; software version (all patient-level undefined; slice-level 0.501--0.527) & series count (slice-level 0.473) & \emph{none in the file} \\
\addlinespace[2pt]
PI-CAI, case level & \emph{slice index absent}; volume depth 0.504 & contributing center 0.547; acquisition year 0.552 & age 0.636; prostate-specific antigen 0.638 \\
\addlinespace[2pt]
LUNA16 & published world $z$; volume depth 0.658 & \emph{none in the file} & \emph{none in the file} \\
\bottomrule
\end{tabular}

\vspace{2pt}
\raggedright\scriptsize
\textbf{Note.}---Class G is the only class present on the hemorrhage benchmark,
which is why the primary result is a statement about geometry and nothing else.
Two arms are not class G alone and the distinction changes how they should be
read. On PI-CAI the depth-limited tree over the four
available variables reaches 0.692 (95\% CI: 0.626, 0.755) at the patient level;
taken one at a time the two clinical variables reach 0.636 and 0.638 and the two
administrative variables reach 0.547 and 0.552, and the fitted tree splits on all
four. The single-variable ordering places the clinical variables above the
administrative ones, but the tree's 0.692 cannot be partitioned between the two
classes from the recorded output, and it is not claimed to be: age and
prostate-specific antigen are established predictors of clinically significant
prostate cancer, so a model using them is not pixel-blind in the same sense as a
model using only stack geometry. On DeepLesion the label
\emph{is} the anatomic region, so relative position is legitimate task
information rather than a shortcut, and that arm is reported as a reference
level and not as a benchmark defect. The distribution-archive bucket on the
fastMRI Prostate T2 arm is the clearest class A result: a field that records only which
download package a subject shipped in ranks patients at 0.596.
\end{table*}

% ---------------------------------------------------------------- Table E2
\begin{table*}[t]
\centering
\footnotesize
\caption{Provenance validation of the RSNA ICH slice ordering and subject identifier, neither of which the official release contains. Each test is
reported against the value the same statistic returns when the structure being
tested is destroyed, so a null result is distinguishable from a test with no
power.}
\label{tab:provenance}
\begin{tabular}{@{}p{0.28\textwidth} p{0.23\textwidth} p{0.23\textwidth} p{0.13\textwidth}@{}}
\toprule
Test & Observed & Value when the structure is destroyed & Separation \\
\midrule
\multicolumn{4}{@{}l}{\textbf{Slice ordering}}\\
\addlinespace[2pt]
Runs of hemorrhage slices per series, \emph{any} label, \num{8377} series & 1.384 & 7.167 expected under random placement; ratio 1.001 (SD 0.002) over 20 randomized orderings & $z=-452$ \\
Same test, subdural & ratio 0.175 & 0.999 (SD 0.003) & $z=-274$ \\
Same test, epidural (rarest, 331 series) & ratio 0.180 & 1.006 (SD 0.010) & $z=-86$ \\
Same test, three remaining subtypes & ratio 0.223, 0.223, 0.246 & --- & all fire \\
Series individually more contiguous than chance & 99.5\% & 50\% by construction & --- \\
Series whose positive slices form one run & 74.8\% & median 7.5 runs expected & --- \\
Sensitivity: ratio after a fraction of slices is relocated at random within its series & 0.193 at 0\%; 0.233 at 5\%; 0.301 at 10\%; 0.444 at 20\%; 0.576 at 30\%; 1.001 at 100\% & pre-set abort threshold 0.50, crossed near 24\% & --- \\
Orientation coherence: mean relative position of positive slices & 0.5239 & 0.4998 (SD 0.0011) if a random half of series are reversed & $z=+22$ \\
\addlinespace[3pt]
\multicolumn{4}{@{}p{0.87\textwidth}}{\textbf{Subject identifier}, tested on \num{1758} subjects holding \num{4564} series; series and studies are 1:1, so every pair below spans two separate examinations}\\
\addlinespace[2pt]
Same-subject series agreeing on reconstruction intercept & 0.977 & 0.515 (SD 0.007) with subject labels permuted & $z=68$ \\
Same-subject series agreeing on hemorrhage present & 0.801 & 0.500 (SD 0.007) & $z=44$ \\
Same, restricted to pairs already matching on reconstruction intercept (\num{5023} pairs) & 0.800 & 0.504 (SD 0.012) & $z=25$ \\
Same-subject series agreeing on series depth & 0.188 & 0.098 (SD 0.002) & $z=40$ \\
\addlinespace[3pt]
\multicolumn{4}{@{}l}{\textbf{Structural reconciliation}}\\
\addlinespace[2pt]
Series whose within-series counter runs 1 to $n$ without a gap & \num{21743} of \num{21744} & --- & --- \\
Rows implied missing by the single gapped series & 1 & reconciles \num{752802} held against \num{752803} published & --- \\
Series present in the labels but absent from the acquisition-parameter file & 0 of \num{21744} & --- & --- \\
\bottomrule
\end{tabular}

\vspace{2pt}
\raggedright\scriptsize
\textbf{Note.}---No image header was read. The $z$ order is inferred from label
contiguity, not observed directly, and the official label file was never obtained,
so the mirror's instance identifiers were not checked row for row against it.
Nothing here identifies which end of the counter is the vertex; the positional
model is invariant to reversing the axis and is fitted per series, so orientation
is not needed. The orientation test is one-sided: asymmetry implies the series
agree on direction, while symmetry would have implied nothing. Every test is
internal-consistency evidence and can only falsify a bad mapping, never certify a
good one---a mapping wrong in a way that preserved within-series label contiguity
and within-subject acquisition homogeneity would pass all of them. The
sensitivity row separates two questions that are easy to conflate. As a
continuous statistic the ordering test is sharp: relocating 5\% of slices moves
the ratio from 0.193 to 0.233, roughly 20 standard deviations of the randomized
null. As a pass/fail gate at the pre-set 0.50 threshold it is blunt: about a
quarter of all slices would have to be misplaced before it refused the file. The
gate certifies against wholesale mis-ordering; it does not certify against a
handful of locally transposed slices, and no claim here depends on that finer
resolution, because the positional model bins relative position into 20 levels.
\end{table*}
